\section{Заключение}

В ходе выполнения данной работы были решены следующие задачи:

\begin{enumerate}
    \item Изучены теоретические основы алгоритмов пузырьковой и быстрой сортировки;
    \item Выведены формулы временной сложности: \texttt{O(n\textsuperscript{2})} для пузырьковой и \texttt{O(n log n)} для быстрой сортировки;
    \item Выполнена программная реализация алгоритмов на языке C\texttt{++};
    \item Проведено экспериментальное исследование производительности;
    \item Выполнен сравнительный анализ полученных результатов.
\end{enumerate}

\textbf{Основные выводы:}

\begin{itemize}
    \item Пузырьковая сортировка из-за квадратичной сложности \texttt{O(n\textsuperscript{2})} неприменима для больших объёмов данных. Она может использоваться только в учебных целях или для очень маленьких массивов.
    
    \item Быстрая сортировка демонстрирует высокую производительность на всех тестовых наборах данных. Её сложность \texttt{O(n log n)} делает её оптимальным выбором для большинства практических задач.
    
    \item Полученные экспериментальные данные полностью подтверждают теоретические оценки сложности. Разница в производительности достигает трёх порядков на массивах размера \texttt{n = 100000}.
\end{itemize}

В ходе выполнения работы были подтверждены теоретические данные, изложенные в трудах Кнута \cite{knuth1998art} и Хоара \cite{hoare1962quicksort}. Справочная информация была взята из документации C\texttt{++} \cite{cppreference}.

\textbf{Перспективы дальнейших исследований:}

\begin{itemize}
    \item Анализ производительности на частично отсортированных и обратно отсортированных данных;
    \item Сравнение с другими эффективными алгоритмами (сортировка слиянием, пирамидальная сортировка, \texttt{Timsort});
    \item Исследование параллельных и многопоточных реализаций алгоритмов сортировки;
    \item Анализ производительности на различных типах данных (строки, структуры, числа с плавающей точкой).
\end{itemize}